\documentclass[10pt,a4paper]{amsart}
\usepackage[margin=19mm]{geometry}
\usepackage{amsmath,amssymb,mathtools}
\usepackage[scr=rsfs,cal=euler]{mathalfa}
\usepackage{xfrac}
\usepackage[unicode,hidelinks,bookmarks=false]{hyperref}
\newcommand{\ie}{i.e.\ }
\newcommand{\cf}{cf.\ }
\newcommand{\resp}{respectively\ }
\newcommand{\Subsection}{\subsection}
\providecommand{\nobreakdash}{\nobreak\hbox{-}}
\setlength{\emergencystretch}{2em}
\multlinegap0pt
\usepackage{bm}
\usepackage{bbm}
\usepackage{dsfont}
\usepackage{xspace}
\usepackage{cancel}
\usepackage{mathtools}
\usepackage{amsthm}
\makeatletter
\@ifundefined{thm}{\newtheorem{thm}{Theorem}}{}
\makeatother
\def\Rb{\overline{{\mathbb R}}_+}
\title[Mass conservation and gelation for the Smoluchowski coagulat]{Mass conservation and gelation for the Smoluchowski coagulation equation: a generalized moment approach}
\author{Masato Kimura, Hisanori Miyata}
\date{Source: arXiv:2506.08017v1 (CC BY 4.0)}
\begin{document}
\maketitle

\noindent\emph{Source and licence.} Mass conservation and gelation for the Smoluchowski coagulation equation: a generalized moment approach. Version arXiv:2506.08017v1 (CC BY 4.0), distributed under the Creative Commons Attribution 4.0 International licence, as recorded in the arXiv licence metadata for this exact version and shown on the abstract page https://arxiv.org/abs/2506.08017v1. Full rights notice: licenses/arxiv-2506.08017-CC-BY-4.0.txt.\par\smallskip

\noindent\textit{Editorial scope.} Counted unit: the Thm whose source label is \texttt{mainth} in the article (source0.tex lines 912--914, source label \texttt{mainth}), delivered together with the article's own complete original proof of it (source0.tex lines 915--993, one \texttt{proof} environment of 79 lines). No mathematics is written, repaired or re-derived: statement and proof are the article's own text line for line. Required context retained uncounted: TMI (lines 555--564); M0inf (lines 638--648); Kr (lines 641--643). Every definition, display and label of those lines is retained unchanged. Omitted from this package and not redistributed: 1--554, 565--637, 644--911, 994--1294. Nothing inside a retained excerpt is omitted or altered: every retained body is delivered exactly as the article's lines are written, so no omission receipt of any kind accompanies this package. Citations: the retained excerpts carry no citation command at all, so no bibliography entry of the article is transcribed and no reference is redistributed. Reference adaptation: none was needed and none was made; every reference inside the delivered unit resolves inside it, so no unresolved reference or citation is produced. Printed slips kept literal: no printed word, symbol, hypothesis or number of the counted statement or of its proof was altered. Layout and class: the article's own class is \texttt{article} with packages including graphicx, amsmath, amsthm, amssymb, enumerate, here, tikz, titlesec, empheq, color, xfrac, bigints, hyperref, subcaption; the wrapper is the lane's pinned \texttt{amsart} editorial wrapper at 10pt on A4, which restates the theorem-like environments the retained text needs and adds the article's own macro definitions that the retained text needs (\texttt{\textbackslash Rb}), with extra wrapper packages (bm, bbm, dsfont, xspace, cancel, mathtools). The delivered numbering is the unit's own. Assets: no retained span contains a figure environment, a graphics-inclusion command, a PDF-page inclusion, a table or a TikZ picture; no image file is copied into this package, the package declares no asset dependency and no image file is redistributed with it. Licence: version arXiv:2506.08017v1 of this article is distributed under the Creative Commons Attribution 4.0 International licence, as recorded in the arXiv licence metadata for this exact version and shown on the abstract page https://arxiv.org/abs/2506.08017v1. A full rights notice is delivered at licenses/arxiv-2506.08017-CC-BY-4.0.txt. Characters: every retained excerpt is pure ASCII, so the delivered engine, XeLaTeX with its default Latin Modern text font, reports no missing character.
\begin{thm}[Truncated generalized moment identity] \label{TMI}
Let $u$ be a weak solution to the SCE on $I=[0,T]$.
Then, for any $b\in C^0(\Rb )$,
the following identity holds:
\begin{align*}
\partial_t m^b (r,t) = \frac{1}{2}\iint_{D(r)} \big(&b(x+y)-b(x)-b(y)\big)K(x,y)u(x,t)u(y,t)\,dxdy \\
&-\int^r_0 b(x)u(x,t)\int^\infty_{r-x} K(x,y)u(y,t)\,dydx\quad(r\in \Rb,~\text{a.e.}~t\in  I)
\end{align*}
where $D(r) \coloneqq \{(x,y)\mid0 \leq x \leq r, 0\leq y \leq r-x\}$.
\end{thm}

\begin{thm}\label{M0inf}
Let $u$ be a weak solution to the SCE on $\Rb$.
If $K(x,y)$ satisfies the following condition:
\begin{align}\label{Kr}
\inf_{x,y\geq r}K(x,y) >0\quad (r>0),
\end{align}
then it follows that
\begin{align}\label{limM0}
\lim_{t\to\infty}M_0(t)=0.
\end{align}
\end{thm}

\begin{thm}\label{mainth}
Suppose that the coagulation kernel $K(x,y)$ satisfies the condition in \eqref{Kr}, and that the weight function $b\in C^0(\Rb;\Rb)$ is non-decreasing on $\Rb$. If condition (A3) is satisfied, then the weak solution $u(x,y)$ of the SCE on $\Rb$ with $u_0\not\equiv 0$ and $M^b(0)<\infty$ undergoes gelation. That is, there exists $T>0$ such that $M_1(t)<M_1(0)$ for $t>T$. 
\end{thm}

\begin{proof}
we begin by noting that $b(x)\ge 0$ and $b(y)\le b(x+y)$,
which implies the inequality:
\begin{align*}
b(x)\ge \frac{1}{2}b(x)\ge \frac{1}{2}(b(x)+b(y)-b(x+y)).
\end{align*}
Applying this in Theorem~\ref{TMI}, we have
\begin{align*}
\partial_t m^b (r,t) &=\frac{1}{2} \int^r_0\int^{r-x}_0 
(b(x+y)-b(x)-b(y))K(x,y)u(x,t)u(y,t)\,dydx \\
&\quad -\int^r_0\int^\infty_{r-x} b(x)K(x,y)u(x,t) u(y,t)\,dydx \\
&\leq \frac{1}{2}\int^r_0\int^{\infty}_0 
(b(x+y)-b(x)-b(y))K(x,y)u(x,t)u(y,t)\,dydx \\
&\leq \frac{1}{2}\int^r_0\int^{\infty}_0 
(-\lambda xy+\mu (x+y+1))u(x,t)u(y,t)\,dydx \\
&=-\frac{\lambda}{2}m_1(r,t)M_1(t)
+\frac{\mu}{2}\left(
m_1(r,t)M_0(t)+m_0(r,t)M_1(t)
+m_0(r,t)M_0(t)\right)\\
&\le 
-\frac{\lambda}{2}m_1(r,t)M_1(t)
+\frac{\mu}{2}M_0(t)(2M_1(t)+M_0(t))\\
&\le 
-\frac{\lambda}{2}m_1(r,t)M_1(t)
+N(t),
\end{align*}
where we define
\begin{align*}
N(t):=\frac{\mu}{2}M_0(t)(2M_1(0)+M_0(0))
\end{align*}
Integrating both sides from $0$ to $t$,
we obtain:
\begin{align}\label{mbtN}
m^b (r,t)-m^b (r,0)
&\le
-\frac{\lambda}{2}\int_{0}^{t}
m_1(r,s)M_1(s)\,ds
+\int_0^tN(s)\,ds.
\end{align}
Since 
\begin{align*}
m^b (r,t)\le m^b (r,0) +\int_0^tN(s)\,ds\le M^b(0)+\int_0^tN(s)\,ds,
\end{align*}
it follows that $M^b(t)< M^b(0)+\int_0^tN(s)\,ds<\infty$.

Now, taking the limit as $r\to\infty$ in \eqref{mbtN} and applying Lebesgue's dominated convergence theorem, we obtain
\begin{align}\label{MMN}
M^b(t)-M^b(0)&\leq
-\frac{\lambda}{2}\int_{0}^{t}
M_1(s)^2\,ds
+\int_0^tN(s)\,ds.
\end{align}

To prove that gelation occurs, we proceed by contradiction. 
Assume that gelation does not occur, i.e., mass is conserved such that $M_1(t)=M_1(0)>0$ for all $t\in\Rb$.
Since $N(t)\to 0$ as $t\to\infty$ by Theorem~\ref{M0inf}, there exists a time $t_*>0$ such that
$N(t)\le \frac{\lambda}{4}M_1(0)^2$ for $t\ge t_*$.
Substituting this into \eqref{MMN} for $t>t_*$, we get:
\begin{align}
M^b(t)&\le M^b(0)-\frac{\lambda}{2}M_1(0)^2t
+\int_0^{t_*}N(s)\,ds+\int_{t_*}^tN(s)\,ds\label{fl}\\
&\le
M^b(0)-\frac{\lambda}{2}M_1(0)^2t
+\int_0^{t_*}N(s)\,ds
+\frac{\lambda}{4}M_1(0)^2(t-t_*)\notag\\
&\le
M^b(0)
+\int_0^{t_*}N(s)\,ds
-\frac{\lambda}{4}M_1(0)^2t\notag
\end{align}
For sufficiently large $t$, specifically when
\begin{align*}
t> \left(M^b(0)
+\int_0^{t_*}N(s)\,ds\right)\frac{4}{\lambda M_1(0)^2}, 
\end{align*}
it follows that $M^b(t)<0$,
which is a contradiction. Thus, mass conservation
does not hold and gelation occurs.
\end{proof}

\end{document}
